% !TeX root = ../colloquium-2.tex

\phantomsection\label{sec_Definition01}
\begin{definition}[Непрерывные на множестве функции. 
      Теорема о нуле непрерывной на отрезке функции.
      Определение ограниченной на множестве функции.]
    \subsection*{Непрерывные на множестве функции.}
    \begin{note}[purple]
        Функция $f$, определённая на множестве $E$, называется непрерывной на $E$,
        если она непрерывна в каждой точке $E$. 
    \end{note}
    \subsection*{Теорема о нуле непрерывной на отрезке функции.}
    \begin{note}[purple]
        Пусть функция $f \in C([a, b])$ и $f(a) \cdot f(b) < 0$. Тогда 
        $\exists c \in (a, b) \; : \; f(c) = 0$. 
    \end{note}
    \subsection*{Определение ограниченной на множестве функции.}
    \begin{note}[purple]
        
    \end{note}
\end{definition}